% check-literals: skip
% Grund: Dieser Abschnitt nennt die Zahl der geprueften Werte und der
% Kettengleichungen. Sie kommen aus den Wachhunden selbst und werden von
% scripts/methode.py in data/methode.tex geschrieben; die Makros stehen im
% Text. Die Ausnahme deckt nur die verbleibenden Gliederungszahlen.
\section*{Methode und Grenzen}
\addcontentsline{toc}{section}{Methode und Grenzen}

\subsection*{Was gepr\"uft ist}

Jeder Zahlenwert dieses Dokuments steht genau einmal in
\texttt{scripts/kennzahlen.py}, zusammen mit seiner Quelle und seiner
gedruckten Seite. Prosa, Tabellen und Abbildungen lesen dieselbe Erkl\"arung;
Text und Tabelle k\"onnen einander deshalb nicht widersprechen.

Vier Skripte erzwingen das, weil Urteilsverm\"ogen bei dieser Menge versagt
und ein Skript nicht:

\begin{itemize}
\item \texttt{pruefe\_seiten.py} sucht jeden Rohwert im Text genau der Seite,
die sein Beleg nennt. Gepr\"uft werden \MethodeWerte{} Angaben; nicht belegt
sind \MethodeFehler{}. Die \MethodeExtern{} Sekund\"arwerte tragen keine
gedruckte Seite und stehen ausserhalb dieser Pr\"ufung.
\item \texttt{reihe\_azn.py} und \texttt{reihe\_rdy.py} pr\"ufen die
Mehrjahresreihen \"uber Kreuz: Jeder Bericht druckt drei Jahre, und die
\"Uberlappung ergibt \MethodeKette{} Gleichungen, die alle aufgehen. Dieser
Test f\"angt genau den Fehler, den die Seitenpr\"ufung nicht sehen kann: eine
Spalte zu weit gegriffen, der Wert steht auf der zitierten Seite und geh\"ort
einem anderen Jahr.
\item \texttt{check\_literals.py} verbietet jede Zahl im Flie{\ss}text, die
nicht aus der Zahlenschicht kommt.
\item \texttt{check\_footnote\_pages.py} und
\texttt{check\_footnote\_groups.py} halten den Fu{\ss}notenapparat sauber.
\end{itemize}

\subsection*{Drei Farben, drei Verantwortlichkeiten}

\begin{itemize}
\item \textbf{Schwarz}: Zahlen, Tabellen, Quellen, Seiten. Daf\"ur steht die
Prim\"arquelle ein. Wer nur die schwarzen Teile liest, erh\"alt die belegte
Lage ohne jede Wertung.
\item \bk{\textbf{Blau}: Einsch\"atzungen im laufenden Text. Daf\"ur steht der
Verfasser ein.}
\item \vd{\textbf{Rot}: die Voten in den Teilen \ref{sec:verdikt-azn} und
\ref{sec:verdikt-rdy}. Eine Empfehlung unter erkl\"arten Annahmen, kein
Befund.}
\end{itemize}

\subsection*{Prim\"ar und sekund\"ar}

Kurse, Wechselkurse, der Branchenindex und der Analystenkonsens stehen in
keinem Gesch\"aftsbericht. Sie fehlen dort nicht, sie geh\"oren dort nicht
hin. Erhoben sind sie aus benannten Sekund\"arquellen mit Abrufdatum und
gepr\"uft, soweit sie sich pr\"ufen lassen: Das aus Kursziel und Aufschlag
zur\"uckgerechnete Bezugsniveau beider Konsenswerte trifft auf den Cent den
Schlusskurs des \Datum{\AznKursTag}. Die Angaben sind also aktuell und nicht
stehengeblieben.

\subsection*{Was nicht gepr\"uft ist}

Die Rechnungen in den Teilen \ref{sec:flow-azn} und \ref{sec:flow-rdy} sind
Modelle. Sie halten den Mittelzufluss konstant und prognostizieren damit
nichts; was sie leisten, ist die Umkehrung der Frage, welches Wachstum
welchen Endwert der heutige Preis voraussetzt. Diese Umkehrung ist nachrechenbar,
ihre Annahmen stehen offen, und das Ergebnis h\"angt an ihnen. Wer die
Annahmen \"andert, \"andert das Ergebnis; deshalb stehen bei beiden Titeln
drei Endwertannahmen nebeneinander statt einer.

Nicht gepr\"uft und nicht abgebildet sind: Steuern des Anlegers,
Quellensteuern, Transaktionskosten, das W\"ahrungsrisiko der Rupie f\"ur
einen Anleger in Dollar oder Euro, und die Frage, wie diese beiden Titel in
ein bestehendes Verm\"ogen passen.

\subsection*{Wann diese Analyse veraltet}

Mit dem n\"achsten Zwischenbericht. F\"ur AstraZeneca ist das der Bericht
zum dritten Quartal 2026, f\"ur Dr.~Reddy's der zum zweiten Quartal des
Gesch\"aftsjahres 2027. Beide Voten h\"angen an Gr\"o{\ss}en, die dort neu
stehen werden. Die Kurse, die H\"urde und der Konsens sind Stichtagswerte des
17.~bzw.\ 18.~September 2026.

\subsection*{Hinweis zur Nutzung von KI-Werkzeugen}

Dieses Dokument ist mit Unterst\"utzung eines Sprachmodells entstanden. Das
Modell hat die Quellen beschafft und seitenweise zerlegt, die Zahlenschicht
und die Pr\"ufskripte geschrieben, die Rechnungen gef\"uhrt und die Prosa
entworfen. Jede Zahl ist maschinell gegen die gedruckte Seite ihrer Quelle
gepr\"uft; kein Wert ist gesch\"atzt oder erg\"anzt worden. Die blau und rot
gesetzten Passagen sind Einsch\"atzungen und Voten und als solche vom
Verfasser zu verantworten, nicht vom Werkzeug.
